\scanpage{120}
--- à tirer au clair à l'occasion !

Revenant au modèle $U$, fournissant un quotient $\mathcal E_U^{\overline L}$ de $\mathcal E_L^{\overline L}$, d'où $N_{x'}\longrightarrow\mathcal E_U^{\overline L}$, induisant $I_{x'}\longrightarrow\mathcal E_U^{\overline L}$,
\deleted{il est clair que si}
qui n'est autre que l'hom \deleted{\unsure{can|$\longleftarrow$}} $\mathcal E_{\widetilde U_x}\unsure{\mathrel{\text{---}}|\longrightarrow}\mathcal E_U$ associé à $\widetilde U_x\longrightarrow U$. \uline{Si} cet hom factorise par $\widetilde S_x$, il est immédiat que l'hom induit $I_{x'}\longrightarrow\mathcal E_U^{\overline L}$ est nul, puisque \inserted{$I_{x'}$} est le noyau de $\mathcal E_{\widetilde U_x}\longrightarrow\mathcal E_{\widetilde S_x}$, et $\mathcal E_{\widetilde U_x}\unsure{\mathrel{\text{---}}|\longrightarrow}\mathcal E_U$ se factorise par cet hom.
\margin{i.e. \uline{si} $V$ a un ``centre sur $U$'' i.e. domine un anneau local de $U$\ldots}
[ \uline{A t-on une réciproque ?} regarder ! ]
\margin{Ça devrait être le cas pour $U$ ``élémentaire'' \unsure{à|et} facteurs \unsure{pas|par} de \uline{\unsure{types $(0,0)$}} \uuline{\unsure{ni|ou}} \uline{$(0,1)$\ldots} \unsure{lié à ceci} \unsure{enfin à l'existence} d'un $K$}

Quand on dispose d'un modèle \uline{projectif} dont $U$ est un ouvert, on peut dire (par \unsure{définition valuative|critère valuatif} de la propreté) que $V$ a nécessairement un centre sur $X$, soit $a$, \deleted{et le centre} \inserted{(\uline{NB} $a$ pas} néc. de codim.1\ldots) et la condition envisagée tantôt (assurant que $I_{x'}\unsure{\mathrel{\text{---}}|\longrightarrow}\mathcal E_U^{\overline L}$ est nul) est que $a\in U$.
\margin{119}
